\section{Renormalizable species counting: a four-tier hierarchy}
\label{sec:hierarchy}

At $t_f\gg t_{\rm fo}$, confinement and the mass gap imply the state is,
up to $O(e^{-m_\pi d})$ corrections, a superposition of well-separated
color-singlet clusters moving ballistically.  The measurement proceeds in
two stages: a geometric \emph{front end} that locates clusters, and a
hierarchy of \emph{species labels} read out on them.

\subsection{Cluster front end and Gauss-law certification}
\label{sec:front-end}

Cluster support is found from the coarse-grained energy density
$\varepsilon_\ell(\mathbf x)$ ($T^{00}$ smeared over
$\ell_{\rm cl}$) --- a heuristic whose scheme dependence is harmless
because all species-defining data comes from the RG-safe operators below;
$\ell_{\rm cl}$-independence of final counts is a reported check.  Regions
$R_i$ are super-threshold domains dilated by a buffer $d$.

A bonus unique to the Hamiltonian formulation: the total color charge in
$R_i$ equals the chromoelectric flux through $\partial R_i$, so measuring
the boundary-link Casimir $\sum_{\partial R_i} E^2$ and finding it
vacuum-consistent \emph{certifies gauge-invariantly that the cluster is a
singlet}, with exponentially small tails
(Appendix~\ref{app:charge-quantization}).

\subsection{Tier 1: conserved-charge counting}
\label{sec:tier1}

On each certified cluster, measure baryon number $B$, strangeness $S$, and
charm $C$ through \emph{exactly conserved point-split lattice vector
currents}, for which $Z\equiv 1$ with no anomalous dimension --- the single
most important technical choice in this construction.  Species logic:
$B=1$ baryon; $B=0$ meson sector; \emph{manifestly exotic by charges
alone}: $(B,C)=(0,+2)$ is at minimum $cc\bar q\bar q$ (the $\Tcc$ class),
$B=2$ a dibaryon, etc.  Hidden-flavor exotics ($\chic$, $P_c$) are
invisible to this tier by construction.

The renormalization statement is a \emph{quantization theorem}
(Appendix~\ref{app:charge-quantization}): on a cluster deep inside $R$
with vacuum near $\partial R$, the distribution of the window-smeared
charge $Q_R^{\ell_w}$ concentrates on an integer with tails
$\le C\,e^{-m_\pi d} + O\big((a/\ell_w)^2\big)$; window smearing over
$\ell_w$ removes the perimeter-divergent contact terms familiar from
charge-fluctuation studies.  Integer labels carry no anomalous dimension;
species numbers inherit RG-invariance from the labels, not from operator
identities.

\subsection{Tier 2: spectral identification}
\label{sec:tier2}

Within $t_{\rm fo}\ll t_f\ll 1/\Gamma$, a narrow exotic \emph{is} a
cluster; its species is identified by rest mass against the independently
computed spectrum in that flavor sector.  Two routes:

\emph{Route S (primary for identification).}  Quantum phase estimation
\cite{Kitaev:1995qy,Choi:2020pdg} on a region-restricted Hamiltonian
$H_R$: the low-lying cluster spectrum is boundary-closure independent up
to $O(e^{-m_\pi d})$, spectral projectors are RG-invariant (in the spirit
of Ref.~\cite{Giusti:2008vb}), the vacuum additive constant is fixed by
running the same QPE on a vacuum region in the same simulation, and
disjoint-region $H_R$'s commute, allowing simultaneous per-shot
identification (Appendix~\ref{app:spectral-projectors}).

\emph{Route T (primary for momenta).}  $P^\mu_R=\int_R T^{0\mu}$ gives
each cluster's four-momentum and hence invariant mass.  The equal-time
renormalized energy-momentum tensor is this construction's largest
technical obligation: we normalize through Ward identities (conservation
across the real-time evolution plus the known total four-momentum of the
prepared state), and defer the 3D small-flow-time program to future work
(Appendix~\ref{app:emt}) \cite{Suzuki:2013gza,Makino:2014taa,
DallaBrida:2020gux}.  Species identification never relies on Route T.

\subsection{Tier 3: asymptotic pair spectra}
\label{sec:tier3}

Defined in Sec.~\ref{sec:resonances}: invariant-mass distributions of
identified-cluster pairs --- the observable experiment actually reports,
and the correct home for broad resonances.

\subsection{Tier 4: flowed coalescence projectors}
\label{sec:tier4}

Instantaneous compact-multiquark projectors built from gradient-flowed
fields \cite{Luscher:2010iy,Makino:2014taa} at fixed physical flow time
$t_{\rm fl}$ have a continuum limit \emph{as a one-parameter labeled
scheme}.  This tier is deliberately scheme-dependent: it is the quantum
image of the classical binding criterion of Sec.~\ref{sec:cluster-id} and
of Lebed's nearest-neighbor rule, and its $t_{\rm fl}$-dependence carries
compact-versus-molecular structure information
(Appendix~\ref{app:flowed}).  Tiers 1--3 count \emph{states}; Tier 4
diagnoses \emph{structure}.
